\section{Discussion and Conclusion}
\label{sec:discussion}\label{sec:conclusion}

\noindent\emph{Discussion.}
Shadow MIA conflates distributional uniqueness with membership: naive
retrain, never trained on the data, scored $0.584$/$0.846$
(N-BaIoT/CASAS). Requiring \emph{all three} MIAs below $\tauMIA$ is
conservative by design, preferring false rejection to false acceptance;
LiRA resists the confound but is $\alpha$-sensitive, so
$(\alpha,\tauMIA)$ are co-calibrated per model class. \sysname accepts
\emph{approximate} unlearning (exact retraining being infeasible at IoT
scale; GDPR Art.~17's ``reasonable measures'' is met by the predicate
plus on-chain audit), leaving formal-erasure and ECDSA-leakage
guarantees (BLS-collapsed in the journal) out of scope; off-chain
training on revoked data is unpreventable but leaves no valid UA, a
detectable omission. On deployment, \sysname discharges GDPR
Arts.~7/17/21, EU AI Act Arts.~15/26, and CCPA \S~1798.105 via the
soulbound NFT, the $\dmax$-bounded liveness and verifiability
predicates, and audit completeness; PoA$^2$ is validator-permissioned
yet publicly readable, while Amoy runs identical contracts
permissionlessly, so the economic model, not the protocol, is the
decision variable.

\smallskip\noindent\emph{Conclusion.}
\sysname is, to our knowledge, the first protocol to bind dynamic
consent, machine unlearning, and on-chain verification in one
chain-enforced pipeline, closing the consent--unlearning--verification
triangle prior work addressed only at the vertices. Across 9-client
N-BaIoT~\cite{meidan2018nbaiot} and 5-client CASAS, approximate
unlearning drives all three MIAs below $\tauMIA\!=\!0.55$ at
$\le\!3\%$ utility loss, with per-subject lifecycle cost $\$0$ on
PureChain versus $\$0.00120$ on Amoy. Beyond the BLS aggregation and
formal-erasure guarantees noted above, future work targets
zero-knowledge eligibility proofs, PoA$^2$-native verifiers, and
cross-chain consent portability.
